\documentclass{article}
\usepackage{ifpdf} 
\usepackage{setspace}
\doublespacing
\title{Midterm Essay}
\author{Mackenzie Norman}
\begin{document}

\textit{``And they buried the hunger artist along with the straw. But in his cage they put a young panther. Even for a person with the dullest mind it was clearly refreshing to see this wild animal throwing itself around in this cage, which had been dreary for such a long time. It lacked nothing. Without thinking about it for any length of time, the guards brought the animal food. It enjoyed the taste and never seemed to miss its freedom. This noble body, equipped with everything necessary, almost to the point of bursting, also appeared to carry freedom around with it. That seem to be located somewhere or other in its teeth, and its joy in living came with such strong passion from its throat that it was not easy for spectators to keep watching. But they controlled themselves, kept pressing around the cage, and had no desire to move on. '' (Kafka, 8)}
\section*{Intro}

\textit{``As if it is ever easy, And these days it is worse, with the poverty of blackness on one side and the weight of womanhood on the other''}(Dangarembga 16)\\ Nervous Conditions: as it concludes, is a story of women. Women who are divided, complex, and despite all living in the same space, women who plot their course through it  very differently. Dangarembga says \textit{``People make individual choices. I think mapping the ground helps in making the choices. Such maps, written in an engaging way, are part of what I perceive some of my responsibility as a novelist to be.''} (Dangarembga 2) In viewing this novel as a map - or rather an atlas - with the characters each as a map. The worldly Nyasha: a rebellious and neurotic youth, our Stephen Daedalus or Quentin Compson and the main character Tambu: an obedient and poor child, with a deep yearning to be free of her colonial shackles.  Finally, Lucia: older and mysterious, with a seemingly magical ability to shape the world to her will. Continuing the metaphor, the roads, the vessels that direct traffic through the map, and how the characters control their fates are cut with food. This essay will inspect how each character cuts her own path, culminating in Lucia's. 


\section*{Nyasha, Oh Starving!}
\textit{``They do not like my language , my English, because it is authentic and my Shona, because it is not!'' }(Dangarembga 196)
\paragraph{}

Nyasha is the character most affected by the titular `nervous conditions'. For Nyasha food is symbolic of so much; locale since she has "tasted" what it feels like to not be in Zimbabwe. This is seen when she returns from England. \textit{``Nyasha tucked into the vegetables with the rest of us. When my mother offered her the sour milk she had asked for, she became very morose. She refused to eat'' }(Dangarembga 52) Nyasha wishes to not be in the physical space that she is and so she rejects the food that is representative of that. 

Food to Nyasha is also a cudgel her father uses to control her. 
\paragraph{}
\noindent
\textit{``I'm not hungry,' Nyasha explained.}\\

\noindent
\textit{"You will eat that food,' commanded the man. "Your mother and I are not killing ourselves working just for you to waste your time playing with boys and then come back and turn up your nose at what we offer. Sit and eat that food. I am telling you. Eat it!' Nyasha took a few mouthfuls.}\\

\noindent
\textit{'She has had enough, Baba,' Maiguru said, but Babamukuru was adamant. He was very upset.}\\

\noindent
\textit{'She must eat her food, all of it. She is always doing this, challenging me. I am her father. If she doesn't want to do what I say, I shall stop providing for her - fees, clothes, food, everything.'}\\


\noindent
\textit{'Nyasha, eat your food,' her mother advised.}

\noindent
\textit{'Christ!' Nyasha breathed and with a shrug picked up her fork and began to eat, slowly at first, then gobbling the food down without a break. The atmosphere lightened with every mouthful she took.''} (Dangarembga 189)
\paragraph{}

While an entire essay could be written on this scene; for this essay, it lays bare how ``the man'' sees not eating the food he provides as a challenge to his authority. In this moment he is not Babamukuru, Nyasha's father, but instead the man or even just men. After Nyasha's last attempt at outward rebellion she almost lost her life. \textit{``In the colonial context the colonist only quits undermining the colonized once the latter have proclaimed loud and clear that white values reign supreme. In the period of decolonization the colonized masses thumb their noses at these very values, shower them with insults and vomit them up.''}(Fanon 8 ) Here she cannot rebel outwardly and instead waits later to vomit it up. The path Nyasha cuts is the most damaging, yet it allows her to retain her dignity. This is solidified when Nyasha later writes to Tambu and says ``[Nyasha] had embarked on a diet "to discipline my body and occupy my mind" ''(Dangarembga 197) Nyashas is characterized as a rebel fighter, and sees her anorexia as training for a final conflict. 

Continuing to inspect food as means for control we know Nyasha has no way to make food on her own. She has neither the skills to grow vegetables that Tambu has nor Lucias apparent ability to conjure it. ``Can you cook books''(Dangarembga 15) is what Tambus father says to her but in the end this applies far better to Nyasha, she is incredibly educated and in the peak of her anorexia she appears to be attempting to live off books. For all Nyashas education, she cannot provide food for herself, so instead she develops anorexia as means to control her destiny. Nyasha's weaponization of food should come as no surprise to us. As Fanon writes \textit{``The very people that had it drummed into them that the only language they understood is force, now decide to express themselves with force. In fact the colonist has always shown them the path they should to liberation.''} (Fanon 42) Nyasha has had food as a mean of control drummed into her for as long as she has been alive. She understands the path to liberation must come through food. 


\section*{ Tambu, The Colonized Intellectual}
Tambus relation to food mirrors her relation to education. Food is a means to escape. We see this literally in the beginning when she is able to pay for her schooling by selling food. When her brother stole her corn, he was stealing her education and as such her ability to control her destiny. She differs from Nyasha in that african food is authentic to her. Compared to Babamakurus family she was the only one who liked Sadza 

\textit{``I was very pleased to see the sadza when it came, although nobody else seemed to care for it. This was embarrassing. There were many things that were embarrassing about that meal: my place looked as though a small and angry child had been fed there; here I was with a spoon in my hand instead of a fork and now Maiguru was dishing sadza on to my plate, sadza that nobody else would eat'' } (Dangarembga 82)


While Nyasha does not find any food authentic, Tambu does. However, she is embarrassed by this and so as Nyasha forces up food, Tambu forces down food. Contrasting Nyasha even more, Tambu takes food as means of control. She eats, even though she is as aware as anyone the weight it carries. 

This embarrassment surrounding food makes clear that Tambu's character is what Fanon would call the ``colonized intellectual'' 
\paragraph{}
\textit{``During this period the intellectual behaves objectively like a vulgar opportunist. His maneuvering, in fact, is still at work. The people would never think of rejecting him or cutting the ground from under his feet. ''} (Fanon 13)
\paragraph{}
As a child - through adolescence Tambu is constantly maneuvering. Her interactions with food represent this. She will eat whatever it takes whether it be potatoes and gravy, or sadza to escape her colonial bonds.

Tambus ability to eat food and assimilate to the patriarchy is buoyed by her ability to create food. Tambu grew her own vegetables from the ground. She is free from the power that Babamukuru wields over Nyasha. But as we will see when we look at Lucia, Tambu is still the colonized intellectual, as she maneuvers towards her own success and only feeds herself instead of helping Nyasha, something that she felt guilty about as evidenced by the passage \textit{``I felt Nyasha needed me but it was true: I had to go to school''}. (Dangarembga 202)  Tambu does not shy away from her original selfishness, as the book begins with  ``I was not sorry when my brother died, Nor am I apologizing for my callousness as you may define it.'' (Dangarembga 1) I think the characterization of Tambu as the ``colonized intellectual'' may seem harsh and in part it is. But Tambu as a map is one of selfishness, and she does not shy away from it in her retrospection. Her selfishness is juxtaposed and exposed by the foil of Lucia, who we will look at next.

\section*{Lucia}
If Tambu and Nyasha lives chart paths around a mountain, Lucia simply climbs it. Despite the topographical lines being so dense they nearly touch, Lucia's path cuts through them. On first glance, this in part can be attributed to her naivete, she truly doesn't know (or at least pretends not to). 

But there is far more to it than that. Firstly Lucia can feed herself, \textit{``Lucia had somehow managed to keep herself plump despite her tribulations''} (Dangarembga 127) and in doing so she is characterized as beautiful, with glowing skin. Something Nhamo only had after he spent a year at the mission. This makes it apparent that Lucia is able to feed herself as well as any man ``Whatever he can do for me, I can do better for myself'' (Dangarembga 145) - even the rich ones. 

Lucia can not only feed herself, but she can feed others. And not just vegetables but meat and milk, something reserved mostly for men. As we look at foods as means of control we see that Babamukuru and other men their ability to provide it as such, but Lucia, existing outside of this binary uses her food to help liberate her people, making herself an embodiment of decolonization \paragraph{}
\textit{``When they returned Lucia cooked the evening meal, making a thick stew out of some meat she had brought with her. "Tambudzai,' she instructed me in my father's presence, 'make sure nobody touches this meat. Nobody besides your mother, who is ill'' }.(Dangarembga 185)
\paragraph{}
 We know Tambus mother is not ill in the literal sense,  she is afflicted with the same ``Nervous Conditions'' that affect so many, but especially Nyasha. This contrasts with Tambu who, despite knowing that Nyasha needs her help (and really her food) returns home with Babamukuru to continue her schooling. Not only is Lucia a foil to the ``colonized intellectual'' But as Fanon writes \textit{``Decolonization ... can, if we want to describe it accurately, be summed up in the well-known words: ``The last shall be first." ''}. (Fanon 2) Lucia is decolonization. She is making the last, first. Going further we see Lucia as the realization that Fanon describes 
\paragraph{}
\textit {``The colonized subject thus discovers that his life, his breathing and his heartbeats are the same as the colonist's. He discovers that the skin of a colonist is not worth more than the "native's." In other words, his world receives a fundamental jolt. The colonized's revolutionary new assurance stems from this. If. in fact, my life is worth as much as the colonist's, his look can no longer strike fear into me or nail me to the spot and his voice can no longer petrify me. I am no longer uneasy in his presence. In reality, to hell with him. Not only does his presence no longer bother me, but I am already preparing to waylay him in such a way that soon he will have no other solution but to flee. The colonial context, as we have said, is characterized by the dichotomy it inflicts on the world.''} (Fanon 10)
\paragraph{}
This is what makes Lucia such a powerful character and powerful `map'. She is the physical embodiment of a realization that she is equal to everything and everyone. As such she exists outside of the binary's that all other characters are forced into. There is a dichotomy forced on the world but she breaks this.  She is somehow feminine and masculine. She is self-serving and unselfish. In an inherent contradiction, she is a contradiction unto it. 




\section*{Conclusion}
The three characters we focus on all use food as a means of control. Tambu, who becomes the ``colonized intellectual'' will eat to connive and to ensure her, and only her `escape'. Even if this means leaving the rebel behind. We find this rebel in Nyasha.  


\paragraph{}
\begin{center}
``
   Mother 

   ... O my son ... an evil and pernicious death.

   Rebel

   Mother, a verdant and sumptuous death.
'' 
\end{center} (Fanon 45)\footnote{ This is from `Les Armes miraculeuses' by Césaire but the only source I can find that is accessible to me is a reddit post with a different translation. So I am currently citing The Wretched of the Earth}
Nyasha is the rebel and she sees in her starvation a verdant and sumptuous death. She has, perhaps shortsightedly, taken up the masters tools and now is trying to dismantle the masters home. This manichean split of characters, which is two maps leading to dead ends is saved by the existence of Lucia. Who is the realization that the artificial dichotomy enforced by colonialism is just that. Artificial. She proves this by existing in both binaries at once. Any `choice' our two protagonists are forced into she avoids. Instead of rejecting food or attempting to use it to serve herself, she uses it to ensure the last are the first. Her skin is described as radiant not only because she is so well fed, but because she is a beacon. One that provides the ``fundamental jolt'' that neither Nyasha nor Tambu can.

\pagebreak
\section*{Works Cited}
    \noindent
    Dangarembga, Tsitsi. Nervous Conditions: A Novel. Seal Press, 2003. 

    \noindent
    Fanon, Frantz. The Wretched of the Earth. Grove Press, 2022. 

    \noindent
    Kafka, Franz. “Franz Kafka Online.” A Hunger Artist by Franz Kafka - Page 8, 1922, www.kafka-online.info/a-hunger-artist-page8.html. 



\end{document}